% Generated draft: requires mathematical and source review.
\documentclass[11pt]{article}
\usepackage[margin=25mm]{geometry}
\usepackage{fontspec}
\setmainfont{Latin Modern Roman}
\usepackage{amsmath,amssymb,mathtools,unicode-math}
\setmathfont{Latin Modern Math}
\usepackage{enumitem,needspace,xurl,hyperref}
\setlength{\parindent}{0pt}
\setlength{\parskip}{5pt}
\setlength{\emergencystretch}{3em}
\newcommand{\N}{\mathbb{N}}
\newcommand{\Z}{\mathbb{Z}}
\newcommand{\Q}{\mathbb{Q}}
\newcommand{\R}{\mathbb{R}}
\newcommand{\C}{\mathbb{C}}
\newcommand{\F}{\mathbb{F}}
\newcommand{\A}{\mathbb{A}}
\newcommand{\PP}{\mathbb{P}}
\newcommand{\E}{\mathbb{E}}
\newcommand{\eps}{\varepsilon}
\newcommand{\dd}{\,\mathrm{d}}
\begin{document}
\section*{UnsolvedMath ID: 701255}
\textbf{Source problem code:} catalogue entry
\subsection{Definitions and mathematical statement}
For a prime \(p\), a finite \(p\)-group is a finite group whose order is \(p^n\) for some integer \(n\geq 0\). For a group \(G\), let \(\operatorname{Aut}(G)\) denote the group, under composition, of all bijective homomorphisms from \(G\) to itself.

The question proposed by D. MacHale is whether there exist an odd prime \(p\), an integer \(n\geq 1\), and a finite group \(G\) such that
\[
|G|=p^n
\qquad\text{and}\qquad
|\operatorname{Aut}(G)|=|G|.
\]
Thus \(G\) must be nontrivial and have odd order. The source also refers to Archive 12.77, whose contents are not supplied.

The supplied registry labels this problem open. Its background and status review report an affirmative claim attributed to an unverified AI solution by A. Kriebel (2026). That report is a source qualification, not an established answer to the question.
\subsection{Short English statement}
Does there exist a nontrivial finite group of odd prime-power order whose automorphism group has exactly the same order? This is D. MacHale's Kourovka Notebook Problem 16.63. The supplied open label and reported unverified affirmative claim are retained without independent verification.
\subsection{Sources}
\begin{itemize}
\item Local UnsolvedMath registry, permanent ID \(701255\), exact source code \texttt{catalogue entry}, titled \emph{Kourovka Notebook Problem 16.63}; supplied complete source record.
\item D. MacHale, problem proposer, 2006. The supplied source attributes Problem 16.63 to \emph{The Kourovka Notebook (Unsolved Problems in Group Theory)}, 16th issue (2006), and cites the 21st issue edition (2026), \texttt{arXiv:1401.0300v48}. The collection description names E. I. Khukhro and V. D. Mazurov as editors.
\item The statement supplies the cross-reference ``See also Archive 12.77.'' No text of that archive entry is supplied.
\item A. Kriebel, \emph{Preprint}, 2026, reported in the supplied background as an ``unverified AI solution'' claiming an affirmative answer: \url{https://kourovkanotebookorg.wordpress.com/wp-content/uploads/2026/09/kourovka_16_63_v1.1.0.pdf}. No preprint title or contents are supplied.
\item Supplied status review dated \texttt{2026-10-16}: canonical status \texttt{open}, evidence classification \texttt{unrefereed\_claim}. Its notes report a \(\#\) marker for an unverified AI solution claim in the October 2026 21st issue, \texttt{arXiv:1401.0300v48}. This assessment is supplied dated curation, not independently verified.
\end{itemize}
\end{document}
